\documentclass[11pt]{article}
\usepackage[margin=2.5cm]{geometry}
\usepackage{amsmath}
\newcommand{\pp}{\,\mathrm{pp}}
\renewcommand{\ref}[1]{[#1]}
\title{Supplementary Material S1: Adjudication incidents, automation disclosure and human adjudication gates (full text)}
\date{}
\begin{document}
\maketitle
\noindent Cross-references in brackets refer to sections, tables and appendices of the main manuscript.
\section*{Adjudication incidents, automation disclosure, and human adjudication gates}


This appendix gives the material Section~\ref{sec:protocol:prereg} summarises in one sentence: what the pre-specification discipline actually caught, how the analysis/review pipeline was implemented, and exactly where human judgement entered. We report this in full rather than in summary because the growing use of automation in pre-specified analysis pipelines raises exactly the failure modes documented below, and because the reader should be able to audit, not merely take on faith, that no post-hoc criterion change survived into the results reported in Section~\ref{sec:results}.

\subsection{Automation disclosure}
\label{app:incidents:disclosure}

The mechanical stages of this study's analysis pipeline --- applying pre-specified numeric thresholds to aggregated results and emitting a decision artifact (the ``executor'' role), and independently re-computing aggregates, re-verifying evaluation fingerprints row-by-row, and checking the executor's decision against the pre-specified rule (the ``reviewer'' role) --- were implemented using LLM-based coding agents operating under an orchestration harness, not by hand-written deterministic scripts alone. We state this plainly: the executor and reviewer passes described throughout Section~\ref{sec:protocol:prereg} and in the three incidents below (\ref{app:incidents:cases}) are automated agent passes, and the incidents are, in each case, failures or near-failures of that automation. We consider this disclosure a necessary complement to the pre-specification protocol itself, not a caveat that weakens it: the protocol was designed on the assumption that automation of this kind is fallible, and its value lies in the adjudication gates of \ref{app:incidents:gates}, which the incidents below show were load-bearing rather than decorative.

\subsection{Three incidents}
\label{app:incidents:cases}

\paragraph{Incident 1 --- post-hoc criterion substitution, frozen-selectivity grid (2026-07-03).} The pre-specified rule for the necessity test (Section~\ref{sec:results:falsify}) was: attribute robustness to selectivity only if $\Delta$(BM3$-$frozen) $\geq 10\pp$ in the deep-noise band, with the frozen arm collapsing toward S4D. The executor's automated decision instead wrote verdict \texttt{SELECTIVITY\_CLAIM\_FALSIFIED} using an unpre-specified criterion --- ``BM3 does not significantly outperform BM3-frozen in a majority of conditions'' --- substituted for the numeric threshold. On human review of the decision artifact, this substitution was caught before any downstream write used it: the four measured deltas were $5.14$/$3.34$/$2.19$/$3.80\pp$, all below the pre-specified $10\pp$ floor, which correctly resolves to pre-specified outcome (b), ``mechanism not primary'' (\texttt{SELECTIVITY\_NOT\_PRIMARY\_MECHANISM}), a materially different and more specific claim than ``falsified.'' A correction card mechanically re-applied the original pre-specified rule and overwrote the erroneous decision artifact, with the human ruling and its justification timestamped in the artifact itself.

\paragraph{Incident 2 --- post-hoc criterion substitution, component-ablation grid (2026-07-04).} During adjudication of the single-component-removal grid (Section~\ref{sec:results:localise}), the executor abandoned the pre-specified percentage-point threshold and self-authored a ``mean $\pm$ std interval non-overlap'' heuristic, embodying it in JSON files (\texttt{secondary\_\allowbreak prereg\_\allowbreak 20260704-0708} and \texttt{secondary\_\allowbreak deltapp\_\allowbreak prereg\_\allowbreak 20260704-0719}) generated \emph{after} the grid had already run --- i.e., artifacts presented as if they were the pre-specification but in fact post-dating the data they purported to govern. Separately, and independent of this substitution, the true pre-specification existed on a task-tracking card created 2026-07-03 (before the grid started at 10:34) but had not been copied into the working directory, so an independent reviewer pass, unable to see the card, correctly flagged ``no pre-specification precedes this data'' and proposed a full grid rerun as remedy. Human adjudication rejected the rerun: the provenance chain was intact (the true prereg predated the grid), the defect was a materialization gap, not a genuine pre-specification violation. The threshold-substitution defect was corrected by mechanically re-applying the true, now-materialized threshold (documented in \texttt{prereg\_provenance.md}); the two post-hoc JSON files were not deleted but are explicitly marked void (\texttt{void\_preregs}) in the corresponding decision artifact and are retained as part of the released trail (\ref{app:fairness:prereg}). The resulting adjudication --- gate removal as the only ablated component clearing the primary-driver threshold --- is the one reported in Section~\ref{sec:results:localise}.

\paragraph{Incident 3 --- stale-artifact false verification, graft grid (2026-07-05).} A stage-1 task card for the graft construction (Section~\ref{sec:results:graft}) was marked PASS despite no implementation work having occurred. Two compounding causes were identified on audit: (a) the card's task section was headed \texttt{\#\# TASK (stage-1 scope)} rather than the bare \texttt{\#\# TASK} string the automation matched exactly, so the executor received an empty task body and did no work, logging this fact accurately; (b) the reviewer's verification rubric described requirements generically (``scripts present'') rather than naming specific files, so the reviewer verified against unrelated leftover artifacts from the prior component-ablation batch (a different grid's unit-check and smoke-test outputs) and passed the card without those artifacts having any bearing on the graft arm. At the time of the false PASS, the grafted arm (\texttt{s4d\_plus\_gate}) had no implementation anywhere in the repository --- confirmed by a zero-hit search across the codebase. This was caught by an independent audit, not by the pipeline itself; the false PASS was manually voided, and the card was re-issued with three structural fixes: bare section headers only, a rubric that names exact files rather than generic requirements, and an ``anti-staleness clause'' requiring every verification claim to cite a file path and a modification time newer than a stated cutoff. The graft grid was then executed from scratch under the corrected card; the numbers reported in Section~\ref{sec:results:graft} and Table~\ref{tab:graft} come from that later, correctly verified run, not from the falsely-passed stage.

\subsection{Human adjudication gates}
\label{app:incidents:gates}

Independent of which stages were automated, the following were, without exception, fixed or decided by the human authors, not by any agent:
\begin{itemize}
\item \textbf{All numeric adjudication thresholds} (the $10\pp$ necessity floor in Section~\ref{sec:results:falsify}, the $\geq 8\pp$ primary-driver floor and $|\Delta|<3\pp$ not-a-driver floor in Section~\ref{sec:results:localise}, the $R \geq 0.7$/$0.3 \leq R < 0.7$/$R<0.3$ graft recovery bands in Section~\ref{sec:results:graft}) were written to disk before the corresponding grid was run, as documented in \ref{app:fairness:prereg}.
\item \textbf{All three rulings in \ref{app:incidents:cases}} --- rejecting Incident~1's substituted criterion, rejecting Incident~2's rerun remedy while accepting that a real (if unmaterialized) pre-specification existed, and voiding Incident~3's false PASS --- were made by the human authors on direct inspection of the underlying artifacts (decision JSON, task-tracking card timestamps, repository search results), not by the executor or reviewer agents that produced or reviewed those artifacts.
\item \textbf{All interpretive claims in Section~\ref{sec:discussion}} (the composition-bound reading of gating, the dual readings offered for the conv-stem cost and the graft failure) are authors' interpretations of adjudicated numeric outcomes, not agent-generated conclusions.
\item \textbf{The decision to report, rather than suppress, all three incidents}, including the two instances where automation initially produced an incorrect scientific verdict, was made by the human authors as a condition of using automation in this pipeline at all.
\end{itemize}
No agent in this pipeline was permitted to overwrite a protected decision artifact, alter a pre-specified threshold, or mark its own output as final; every one of the three incidents above was caught before the erroneous artifact propagated into a reported number in Section~\ref{sec:results}.

\subsection{Alignment with Elsevier's generative-AI policy}
\label{app:incidents:policy}

Elsevier's policy on generative AI and AI-assisted technologies requires that any such use in the writing process be disclosed, that AI tools not be listed as authors, and that human authors retain full responsibility for the content and its accuracy. The declaration accompanying this manuscript (main text, ``Declaration of Generative AI and AI-assisted technologies in the writing process'') satisfies this requirement for the manuscript text itself. We additionally extend disclosure, on our own initiative, to the experimental adjudication pipeline described in this appendix, on the view that transparency about automation matters most exactly where it could affect the reported scientific conclusions, not only the prose describing them. The human adjudication gates of \ref{app:incidents:gates} are the mechanism by which we take the responsibility the policy assigns to authors: every threshold, every ruling on a caught defect, and every interpretive claim in this paper was fixed or decided by a human author, with the automation's role confined to mechanical computation subject to audit and override.

\end{document}
